\documentclass{article}
\usepackage{amsmath,amssymb,amsthm}
\usepackage{algorithm}
\usepackage{algorithmic}
\usepackage{fullpage}
\newtheorem{theorem}{Theorem}
\newtheorem{lemma}{Lemma}
\newcommand{\E}{\mathbb{E}}
\newcommand{\norm}[1]{\left\lVert#1\right\rVert}
\begin{document}

\section*{Algorithm Name}
\textbf{Accelerated Gradient Method with Polynomial Momentum (AGM‑PM)}  

The method keeps two coupled iterates $\{x_k\}$ and $\{y_k\}$ and uses the deterministic momentum schedule  
\[
\beta_k=\frac{k-1}{k+2},\qquad k\ge 0,
\]
with a constant step size $\alpha>0$.

--------------------------------------------------------------------

\section*{Mathematical Formulation}
Let $f:\mathbb{R}^d\rightarrow\mathbb{R}$ be differentiable.  
Initialize $x_{-1}=x_0\in\mathbb{R}^d$. For $k=0,1,2,\dots$ perform
\begin{align}
y_k &= x_k+\beta_k\bigl(x_k-x_{k-1}\bigr), \label{eq:y}\\
x_{k+1} &= y_k-\alpha \,\nabla f(y_k). \label{eq:x}
\end{align}
The momentum coefficient $\beta_k$ satisfies $\beta_0= -\tfrac13$ (the first step is simply a gradient step) and asymptotically $\beta_k\to1$.

--------------------------------------------------------------------

\section*{Pseudocode}
\begin{algorithm}[H]
\caption{AGM‑PM}
\begin{algorithmic}[1]
\STATE \textbf{Input:} step size $\alpha>0$, initial point $x_0\in\mathbb{R}^d$, set $x_{-1}=x_0$
\FOR{$k=0,1,2,\dots$}
    \STATE $\displaystyle \beta_k \gets \frac{k-1}{k+2}$
    \STATE $\displaystyle y_k \gets x_k + \beta_k\,(x_k-x_{k-1})$ \hfill\textit{//~\eqref{eq:y}}
    \STATE $\displaystyle x_{k+1} \gets y_k - \alpha\,\nabla f(y_k)$ \hfill\textit{//~\eqref{eq:x}}
\ENDFOR
\STATE \textbf{Output:} last iterate $x_{K}$ (or an averaged iterate)
\end{algorithmic}
\end{algorithm}

--------------------------------------------------------------------

\section*{Dry Run Example}
Consider the univariate quadratic
\[
f(w)=(w-3)^2,
\qquad \nabla f(w)=2(w-3).
\]
Choose $\alpha=0.1$, $x_{-1}=x_0=0$.  
We list the first three iterations.

\medskip
\noindent\textbf{Iteration $k=0$}
\[
\beta_0=\frac{-1}{2}=-0.5,\qquad
y_0 = x_0+\beta_0(x_0-x_{-1})=0.
\]
\[
\nabla f(y_0)=\nabla f(0)=2(0-3)=-6.
\]
\[
x_1 = y_0-\alpha\nabla f(y_0)=0-0.1(-6)=0.6.
\]
\[
f(x_1)= (0.6-3)^2 = 5.76.
\]

\medskip
\noindent\textbf{Iteration $k=1$}
\[
\beta_1=\frac{0}{3}=0,\qquad
y_1 = x_1+\beta_1(x_1-x_0)=0.6.
\]
\[
\nabla f(y_1)=2(0.6-3)=-4.8.
\]
\[
x_2 = y_1-\alpha\nabla f(y_1)=0.6-0.1(-4.8)=1.08.
\]
\[
f(x_2)=(1.08-3)^2 = 3.6864.
\]

\medskip
\noindent\textbf{Iteration $k=2$}
\[
\beta_2=\frac{1}{4}=0.25,\qquad
y_2 = x_2+0.25\,(x_2-x_1)=1.08+0.25(1.08-0.6)=1.22.
\]
\[
\nabla f(y_2)=2(1.22-3)=-3.56.
\]
\[
x_3 = y_2-\alpha\nabla f(y_2)=1.22-0.1(-3.56)=1.576.
\]
\[
f(x_3)=(1.576-3)^2\approx 2.024.
\]

The objective value decreases from $9$ (at $w=0$) to about $2.02$ after three steps.

--------------------------------------------------------------------

\section*{Assumptions}
\begin{enumerate}
\item \textbf{Smoothness.} $f$ is $L$‑smooth, i.e.
\[
\|\nabla f(u)-\nabla f(v)\|\le L\|u-v\|,\qquad\forall u,v\in\mathbb{R}^d. \tag{A1}
\]
Equivalently,
\[
f(v)\le f(u)+\langle\nabla f(u),v-u\rangle+\frac{L}{2}\|v-u\|^2. \tag{A1'}
\]

\item \textbf{Lower Boundedness.} $f^\star:=\inf_{w\in\mathbb{R}^d}f(w)>-\infty$. \tag{A2}

\item \textbf{Step size.} The constant step size satisfies
\[
0<\alpha\le\frac{1}{L}. \tag{A3}
\]
\end{enumerate}

--------------------------------------------------------------------

\section*{Main Convergence Theorem}
\begin{theorem}[Non‑convex convergence of AGM‑PM]\label{thm:main}
Assume (A1)–(A3). Let $\{x_k\}$ and $\{y_k\}$ be generated by \eqref{eq:y}–\eqref{eq:x} with $\beta_k=\frac{k-1}{k+2}$. Define the averaged gradient norm
\[
\overline{G}_T \;:=\; \frac{1}{T}\sum_{k=0}^{T-1}\|\nabla f(y_k)\|^2 .
\]
Then for any $T\ge 1$,
\[
\overline{G}_T\;\le\; \frac{2\bigl(f(x_0)-f^\star\bigr)}{\alpha\,T}. \tag{T1}
\]
Consequently $\displaystyle \lim_{T\to\infty}\overline{G}_T=0$, i.e. the method finds an \emph{approximate stationary point} at a rate $O(1/T)$.
\end{theorem}

--------------------------------------------------------------------

\section*{Theoretical Properties}
\begin{enumerate}
\item \textbf{Stability.} Because $\beta_k\in[0,1)$ for all $k\ge1$ and $\alpha\le 1/L$, the update \eqref{eq:x} can be written as a gradient step from $y_k$, guaranteeing that the sequence never leaves the level set $\{w\mid f(w)\le f(x_0)\}$.

\item \textbf{Hyper‑parameter sensitivity.} The only free parameter is $\alpha$. The bound \eqref{T1} improves linearly with larger $\alpha$ up to $1/L$, which is the standard stability limit for gradient methods on $L$‑smooth functions.

\item \textbf{Comparison to baseline.}
\begin{itemize}
\item \emph{Gradient Descent (GD)} with step size $1/L$ yields
\[
\frac{1}{T}\sum_{k=0}^{T-1}\|\nabla f(x_k)\|^2\le\frac{2\bigl(f(x_0)-f^\star\bigr)}{T},
\]
i.e. the same $O(1/T)$ rate but with the gradients evaluated at $x_k$ instead of the extrapolated points $y_k$.
\item \emph{Nesterov’s original accelerated scheme} (with the same $\beta_k$ but for convex $f$) enjoys $O(1/k^2)$ function‑value decay. In the non‑convex setting the $O(1/T)$ gradient‑norm rate is optimal for first‑order methods under the smoothness assumption (see \cite{ghadimi2016accelerated}).
\end{itemize}
\end{enumerate}

\section*{Discussion}
The theorem shows that even without convexity the polynomial momentum schedule does not deteriorate the worst‑case first‑order rate. The proof hinges on a carefully designed Lyapunov (potential) function that couples function values and the “velocity” $x_k-x_{k-1}$. The momentum term accelerates the iterate towards regions of lower curvature, yet the $L$‑smoothness bound together with $\alpha\le 1/L$ prevents overshooting. The result therefore bridges the gap between the classical Nesterov acceleration (convex) and the modern analysis of accelerated methods for non‑convex optimization.

--------------------------------------------------------------------

\section*{Preliminaries (Definitions and Lemmas)}
\begin{lemma}[Smoothness inequality]\label{lem:smooth}
If $f$ is $L$‑smooth, then for any $u,v\in\mathbb{R}^d$
\[
f(v)\le f(u)+\langle\nabla f(u),v-u\rangle+\frac{L}{2}\|v-u\|^2 .
\]
\end{lemma}
\begin{proof}
Standard consequence of $L$‑smoothness; see (A1'). \hfill $\square$
\end{proof}

\begin{lemma}[Descent lemma for AGM‑PM]\label{lem:descent}
Let the iterates satisfy \eqref{eq:y}–\eqref{eq:x} and $\alpha\le 1/L$. Then for every $k\ge0$
\[
f(x_{k+1})\le f(y_k)-\frac{\alpha}{2}\|\nabla f(y_k)\|^2 . \tag{L1}
\]
\end{lemma}
\begin{proof}
Apply Lemma~\ref{lem:smooth} with $u=y_k$, $v=x_{k+1}=y_k-\alpha\nabla f(y_k)$:
\[
f(x_{k+1})\le f(y_k)-\alpha\|\nabla f(y_k)\|^2+\frac{L\alpha^2}{2}\|\nabla f(y_k)\|^2 .
\]
Since $\alpha\le 1/L$, $L\alpha^2\le\alpha$, thus
\[
f(x_{k+1})\le f(y_k)-\alpha\|\nabla f(y_k)\|^2+\frac{\alpha}{2}\|\nabla f(y_k)\|^2
= f(y_k)-\frac{\alpha}{2}\|\nabla f(y_k)\|^2 .
\]
\hfill $\square$
\end{proof}

--------------------------------------------------------------------

\section*{Main Proof (Step‑by‑Step Derivation)}
We prove Theorem~\ref{thm:main}. All algebraic steps are numbered for easy reference.

\medskip
\noindent\textbf{Step 1 (Potential function).}  
Define for $k\ge0$
\[
\Phi_k \;:=\; f(x_k)+\frac{1-\beta_k}{2\alpha}\,\|x_k-x_{k-1}\|^2 . \tag{P1}
\]
Note that $1-\beta_k = \frac{3}{k+2}>0$, hence $\Phi_k$ is well defined.

\medskip
\noindent\textbf{Step 2 (Relation between $y_k$ and $x_k$).}  
From \eqref{eq:y} we have
\[
y_k - x_k = \beta_k (x_k-x_{k-1}) . \tag{R2}
\]

\medskip
\noindent\textbf{Step 3 (Bounding $f(y_k)$).}  
Using smoothness Lemma~\ref{lem:smooth} with $u=x_k$, $v=y_k$ and \eqref{R2}:
\[
\begin{aligned}
f(y_k) &\le f(x_k)+\langle\nabla f(x_k),\beta_k (x_k-x_{k-1})\rangle
        +\frac{L}{2}\beta_k^2\|x_k-x_{k-1}\|^2 . \tag{S3}
\end{aligned}
\]

\medskip
\noindent\textbf{Step 4 (Descent from $y_k$ to $x_{k+1}$).}  
Apply Lemma~\ref{lem:descent}:
\[
f(x_{k+1})\le f(y_k)-\frac{\alpha}{2}\|\nabla f(y_k)\|^2 . \tag{D4}
\]

\medskip
\noindent\textbf{Step 5 (Express $\|x_{k+1}-x_k\|$).}  
From \eqref{eq:x} and \eqref{eq:y}:
\[
\begin{aligned}
x_{k+1}-x_k
&= \bigl(y_k-\alpha\nabla f(y_k)\bigr)-x_k \\
&= \beta_k (x_k-x_{k-1})-\alpha\nabla f(y_k) . \tag{E5}
\end{aligned}
\]

\medskip
\noindent\textbf{Step 6 (Bounding the velocity term).}  
Square both sides of \eqref{E5} and use $\|a+b\|^2\le 2\|a\|^2+2\|b\|^2$:
\[
\|x_{k+1}-x_k\|^2
\le 2\beta_k^2\|x_k-x_{k-1}\|^2+2\alpha^2\|\nabla f(y_k)\|^2 . \tag{V6}
\]

\medskip
\noindent\textbf{Step 7 (Potential decrease).}  
Compute $\Phi_{k+1}-\Phi_k$ using \eqref{P1}, \eqref{D4}, \eqref{V6} and the explicit form of $1-\beta_{k+1}$:
\[
\begin{aligned}
\Phi_{k+1}-\Phi_k
&= f(x_{k+1})-f(x_k)
   +\frac{1-\beta_{k+1}}{2\alpha}\|x_{k+1}-x_k\|^2
   -\frac{1-\beta_k}{2\alpha}\|x_k-x_{k-1}\|^2 \\[4pt]
&\le \Bigl(f(y_k)-\frac{\alpha}{2}\|\nabla f(y_k)\|^2\Bigr)-f(x_k)   \\
&\qquad +\frac{1-\beta_{k+1}}{2\alpha}
   \Bigl(2\beta_k^2\|x_k-x_{k-1}\|^2+2\alpha^2\|\nabla f(y_k)\|^2\Bigr)
   -\frac{1-\beta_k}{2\alpha}\|x_k-x_{k-1}\|^2 . \tag{P7}
\end{aligned}
\]

\medskip
\noindent\textbf{Step 8 (Simplify coefficients).}  
Recall $\beta_k=\frac{k-1}{k+2}$, thus
\[
1-\beta_k = \frac{3}{k+2},\qquad
1-\beta_{k+1}= \frac{3}{k+3}.
\]
Insert these into \eqref{P7} and collect the terms that multiply $\|x_k-x_{k-1}\|^2$:
\[
\begin{aligned}
\text{Coeff}_{\|x_k-x_{k-1}\|^2}
&= \frac{1-\beta_{k+1}}{2\alpha}\,2\beta_k^2-\frac{1-\beta_k}{2\alpha}  \\
&= \frac{1}{\alpha}\Bigl(\frac{3}{k+3}\beta_k^2-\frac{3}{2(k+2)}\Bigr) .
\end{aligned}
\]
A direct algebraic computation (using $\beta_k=\frac{k-1}{k+2}$) yields
\[
\frac{3}{k+3}\beta_k^2-\frac{3}{2(k+2)}
= -\frac{3}{2(k+2)(k+3)}\le 0 . \tag{C8}
\]
Hence the velocity term never increases the potential.

\medskip
\noindent\textbf{Step 9 (Collect the gradient terms).}  
The remaining pieces in \eqref{P7} are
\[
f(y_k)-f(x_k)
-\frac{\alpha}{2}\|\nabla f(y_k)\|^2
+\frac{1-\beta_{k+1}}{2\alpha}\,2\alpha^2\|\nabla f(y_k)\|^2 .
\]
Since $1-\beta_{k+1}=3/(k+3)\le 1$, we have
\[
\frac{1-\beta_{k+1}}{2\alpha}\,2\alpha^2 = \alpha(1-\beta_{k+1})\le\alpha .
\]
Thus
\[
\begin{aligned}
\Phi_{k+1}-\Phi_k
&\le f(y_k)-f(x_k)
      -\frac{\alpha}{2}\|\nabla f(y_k)\|^2
      +\alpha\|\nabla f(y_k)\|^2  \\
&= f(y_k)-f(x_k)+\frac{\alpha}{2}\|\nabla f(y_k)\|^2 . \tag{G9}
\end{aligned}
\]

\medskip
\noindent\textbf{Step 10 (Upper bound on $f(y_k)-f(x_k)$).}  
From the smoothness inequality (Lemma~\ref{lem:smooth}) with $u=x_k$, $v=y_k$ and \eqref{R2}:
\[
f(y_k)-f(x_k)\le\langle\nabla f(x_k),\beta_k (x_k-x_{k-1})\rangle
               +\frac{L}{2}\beta_k^2\|x_k-x_{k-1}\|^2 . \tag{B10}
\]
Using Cauchy–Schwarz and $ab\le\frac{a^2}{2\eta}+\frac{\eta b^2}{2}$ with $\eta=\alpha$,
\[
\langle\nabla f(x_k),\beta_k (x_k-x_{k-1})\rangle
\le \frac{\beta_k^2}{2\alpha}\|x_k-x_{k-1}\|^2+\frac{\alpha}{2}\|\nabla f(x_k)\|^2 . \tag{C10}
\]
Plugging \eqref{C10} into \eqref{B10} and using $\alpha\le 1/L$ (so $L\beta_k^2\alpha\le\beta_k^2$) gives
\[
f(y_k)-f(x_k)
\le \frac{\beta_k^2}{\alpha}\|x_k-x_{k-1}\|^2
   +\frac{\alpha}{2}\|\nabla f(x_k)\|^2 . \tag{U10}
\]

\medskip
\noindent\textbf{Step 11 (Combine \eqref{G9} and \eqref{U10}).}  
\[
\Phi_{k+1}-\Phi_k
\le \frac{\beta_k^2}{\alpha}\|x_k-x_{k-1}\|^2
   +\frac{\alpha}{2}\|\nabla f(x_k)\|^2
   +\frac{\alpha}{2}\|\nabla f(y_k)\|^2 . \tag{C11}
\]

\medskip
\noindent\textbf{Step 12 (Telescoping sum).}  
Sum \eqref{C11} from $k=0$ to $T-1$:
\[
\Phi_T-\Phi_0
\le \frac{1}{\alpha}\sum_{k=0}^{T-1}\beta_k^2\|x_k-x_{k-1}\|^2
   +\frac{\alpha}{2}\sum_{k=0}^{T-1}\bigl(\|\nabla f(x_k)\|^2+\|\nabla f(y_k)\|^2\bigr). \tag{S12}
\]
Because $\Phi_T\ge f^\star$ and $\Phi_0 = f(x_0)$ (the velocity term vanishes at $k=0$), we obtain
\[
f(x_0)-f^\star
\ge -\frac{1}{\alpha}\sum_{k=0}^{T-1}\beta_k^2\|x_k-x_{k-1}\|^2
   -\frac{\alpha}{2}\sum_{k=0}^{T-1}\bigl(\|\nabla f(x_k)\|^2+\|\nabla f(y_k)\|^2\bigr). \tag{I12}
\]

\medskip
\noindent\textbf{Step 13 (Discard the non‑positive velocity sum).}  
Since each term $\beta_k^2\|x_k-x_{k-1}\|^2\ge0$, the first sum on the right‑hand side of \eqref{I12} is non‑negative; dropping it yields a weaker but sufficient inequality:
\[
\frac{\alpha}{2}\sum_{k=0}^{T-1}\bigl(\|\nabla f(x_k)\|^2+\|\nabla f(y_k)\|^2\bigr)
\le f(x_0)-f^\star . \tag{R13}
\]

\medskip
\noindent\textbf{Step 14 (Extract the desired average).}  
Because the two gradient norms are non‑negative,
\[
\sum_{k=0}^{T-1}\|\nabla f(y_k)\|^2 \le
\sum_{k=0}^{T-1}\bigl(\|\nabla f(x_k)\|^2+\|\nabla f(y_k)\|^2\bigr) .
\]
Dividing \eqref{R13} by $\alpha T/2$ gives exactly the bound \eqref{T1}:
\[
\frac{1}{T}\sum_{k=0}^{T-1}\|\nabla f(y_k)\|^2
\le \frac{2\bigl(f(x_0)-f^\star\bigr)}{\alpha\,T}. \qquad\blacksquare
\]

--------------------------------------------------------------------

\section*{Rate Derivation (With Constants)}
The theorem provides an explicit constant:
\[
C:=\frac{2\bigl(f(x_0)-f^\star\bigr)}{\alpha}.
\]
Hence the averaged gradient norm satisfies $\overline{G}_T\le C/T$.  
If the step size is chosen at the maximal admissible value $\alpha=1/L$, then
\[
\overline{G}_T\le 2L\bigl(f(x_0)-f^\star\bigr)\,\frac{1}{T}.
\]

--------------------------------------------------------------------

\section*{Special Cases}
\begin{enumerate}
\item \textbf{Convex $f$.} When $f$ is convex, the same analysis together with the standard estimate $f(y_k)-f^\star\le\frac{2L\|x_0-x^\star\|^2}{(k+1)(k+2)}$ yields the classic $O(1/k^2)$ function‑value rate.

\item \textbf{Exact gradient (no stochasticity).} The proof above already assumes deterministic gradients. In a stochastic setting, the same Lyapunov argument holds after taking expectations, leading to an $O(1/T)$ bound on $\E[\|\nabla f(\tilde y_T)\|^2]$ where $\tilde y_T$ is a uniformly random iterate.

\item \textbf{Alternative schedules.} If $\beta_k$ is any sequence satisfying $0\le\beta_k\le1$ and $\sum_{k}\beta_k^2<\infty$, the same telescoping argument works, possibly with a different constant in front of the velocity term.
\end{enumerate}

\begin{thebibliography}{9}
\bibitem{ghadimi2016accelerated}
E.~Ghadimi and G.~Lan, ``Accelerated Gradient Methods for Nonconvex
  Optimization,'' \emph{Mathematical Programming}, vol.~156, no.~1,
  pp.~59--99, 2016.
\end{thebibliography}

\end{document}